% ── CHAPTER 3: THE WRITINGS: MAHABBAT, THE TWO WINGS, AND THE NATURE OF REALITY (~8,000 words) ──
\chapter[The Writings: \textit{Mahabbat} and the Two Wings]{The Writings: \textit{Mahabbat}, the Two Wings, and the Nature of Reality}
\label{ch:writings-mahabbat}

\epigraph{All created things are expressions of the affinity and cohesion of
  elementary substances, and nonexistence is the absence of their attraction
  and agreement.}{\AbdulBaha, \textit{The Promulgation of Universal Peace}}

% ─────────────────────────────────────────────────────────────────────────────
%  OPENING
% ─────────────────────────────────────────────────────────────────────────────

\firstword{O}{f} the three traditions this book brings into conversation, the one most
readers will know least about is the \bahai Faith.
Quantum mechanics has a cultural presence far beyond laboratories.
Alfred North Whitehead, while not a household name, belongs to the
recognizable terrain of Western philosophy.
But the \bahai Faith --- founded in Persia in the mid-nineteenth century, now
practiced, on some estimates, by eight million people in virtually every
country on earth --- remains largely unfamiliar to Western readers despite more than a century and
a half of intellectual and social activity.

That unfamiliarity has a cost.
When I set the \bahai writings beside quantum mechanics and process philosophy
in the chapters that follow, I want the reader to encounter them as
what they are: a body of philosophical and spiritual reflection of unusual
depth, precision, and scope --- not as curiosities or as exotic reinforcement
for ideas that came from elsewhere.
The convergence I document has intellectual force only if the reader
can see that the \bahai claims are genuinely independent of the other two
traditions, and that independence requires some acquaintance with what
the writings actually say and when they said it.

This chapter introduces the \bahai tradition to a reader who may be
encountering it for the first time.
I do not aim to be comprehensive.
I aim to do two things.
First, to give a sense of the historical and intellectual context in which
these writings emerged --- who wrote them, under what circumstances, and
with what range of concerns.
Second, to read three specific passages closely enough that their
philosophical content becomes clear.
These three passages form the textual foundation of everything that follows.
Along the way I say what kind of texts they are, how much weight their words
can bear, and what the strongest objection to my reading is.

% ─────────────────────────────────────────────────────────────────────────────
\section{The Historical and Intellectual Context}
\label{sec:ch3-context}
% ─────────────────────────────────────────────────────────────────────────────

Mírzá Husayn-`Alí Núrí was born in Tehran in 1817 into a family of
Persian nobility.
His father was a government minister; he was educated in the manner of his
class, which meant tutoring in Persian literature, calligraphy, and the
conventions of courtly life, but not in the formal curriculum of
Islamic theology or philosophy.
He would later write in both Persian and Arabic with a command that struck
readers trained in those traditions as extraordinary, and sometimes as
inexplicable.
He is known to history by the title he adopted: \Bahaullah, an Arabic phrase
meaning ``the Glory of God.''

By the time he was thirty, \Bahaullah had become a follower of the Bábí
movement, a religious reform movement in Shí`ah Islam that had been founded
in 1844 by a young merchant from Shiraz known as the Báb.
The Bábí movement attracted fierce opposition from both civil and religious
authorities.
The Báb was executed by firing squad in 1850 --- an event witnessed by
thousands in the public square of Tabriz.
\Bahaullah survived, but not without consequences: he was imprisoned in the
notorious black pit of Tehran known as the Síyáh-Chál, then exiled from Persia
to Baghdad, then to Constantinople, then to Adrianople, and finally, in 1868,
to the prison city of `Akká in Ottoman Palestine --- now the city of Acre in
northern Israel.
He spent the remaining years of his life, until his death in 1892, in the
vicinity of `Akká.

He wrote throughout all of it.

The body of writing that emerged from these decades of imprisonment and exile
is enormous by any standard.
Scholars estimate that \Bahaullah's authenticated writings exceed the combined
length of the Bible and the Quran.
The range is as wide: it includes mystical poetry, letters of
governance addressed to kings and religious leaders, detailed prescriptions
for social organization, metaphysical treatises, prayers, and works of
philosophical reflection.
The mode varies with the audience and the occasion, but a single animating
vision runs through all of it: that reality is one, that the human experience
of love is not an accident or an evolutionary convenience but a fundamental
structural feature of the universe, and that humanity is in the early stages
of a collective coming-of-age whose full expression lies ahead.

Their standing within the tradition rests on the claim \Bahaullah made for
them.
The Báb had taught that his own mission prepared the way for a greater
figure still to come.
In 1863, in a garden outside Baghdad on the eve of his removal to
Constantinople, \Bahaullah announced to a few companions that he was that
figure.
\bahai belief regards him as a Manifestation of God, the most recent in a line
of messengers that includes Abraham, Moses, the Buddha, Jesus and Muhammad,
and regards his writings as revelation.
This book does not ask the reader to share that belief, but it explains the
weight his words carry within the tradition.

\Bahaullah's eldest son, `Abbás Effendí, was born in Tehran in 1844 on the
same night that the Bábí movement was founded --- a coincidence that \bahai
tradition has treated as significant.
`Abbás Effendí grew up in the company of his father's exile, shared the
imprisonment in `Akká, and was appointed by \Bahaullah as the authorized
interpreter of his writings.
He is known to history by the title his father gave him: \AbdulBaha,
meaning ``Servant of Bahá.''

\AbdulBaha's role in the tradition is philosophically important and
often misunderstood.
He was not a second founder.
He was not a prophet in his own right.
He was the person whom \Bahaullah designated as the one whose
interpretation of the writings carried authoritative weight --- the one
who could, in the decades that followed, explain what the writings meant
to audiences whose languages, cultures, and intellectual contexts were
radically different from those in which the writings had been composed.
He held that role from his father's death in 1892 until his own, in Haifa,
in 1921.

\AbdulBaha took that role seriously in a way that had lasting intellectual
consequences.
After being released from prison in 1908, following the constitutional
revolution that ended the Ottoman Sultan's absolute rule, he traveled
extensively: to Egypt, then to Europe, then to the United States and Canada.
Between 1911 and 1913, he delivered hundreds of talks in Paris, London, and
cities across North America, and many were written down by those who heard
them and later published.
The two collections that figure most prominently in this book are
\textit{Paris Talks}~\citeyearpar{ParisTalks} and
\textit{The Promulgation of Universal Peace}~\citeyearpar{PUP}, which
collects talks given during the North American journey.

The range of his audiences shaped what he said.
\textit{The Promulgation of Universal Peace} alone records talks at Columbia,
Stanford and Howard universities, at the Bowery Mission in New York, at the
annual conference of the NAACP in Chicago, at a synagogue in San Francisco, at
a woman's suffrage meeting, and before Socialists and labor leaders in
Montreal~\citep{PUP}.
Much of what he told them concerned the social teachings of his father: the
oneness of humanity, the equality of women and men, the harmony of science
and religion, the abolition of prejudice, universal peace.
But he grounded those teachings, again and again, in claims about the nature
of reality.
The result is a body of philosophical reflection, delivered in the idiom
of popular address but reasoning at a level that rewards careful analysis,
on topics ranging from the nature of the soul to the structure of knowledge
to the relationship between religion and science.

\subsection*{Writings, Talks, and the Weight of Words}

Because this book leans on exact phrases, the reader should know what kind of
text each phrase comes from and how much weight it can bear.
The tradition itself draws the relevant lines.

At the center stand the writings of \Bahaullah.
The theme of this book is there in his own voice: ``I loved thy creation,
hence I created thee''~\citep[Arabic no.~4]{HiddenWords}.
The preface to \textit{Selections from the Writings of \AbdulBaha}, compiled
by the Research Department of the Universal House of Justice, sorts
\AbdulBaha's legacy into three kinds: ``His written works,'' ``the many
compilations of His recorded utterances,'' and ``His
correspondence''~\citep[Preface]{SelectionsAbdul}.
The written works and the letters are texts he wrote or approved.
The recorded utterances are what others heard him say and wrote down.

\textit{The Promulgation of Universal Peace} is a book of recorded utterances.
\AbdulBaha spoke in Persian, interpreters put his words into English as he
went, and what survives, for most talks, is a listener's notes of the English.
In the text I have used, 138 of the 140 talks carry a line naming their
source.
Most read ``Notes by'' followed by a name: Joseph H. Hannen, Edna McKinney,
Esther Foster, Howard MacNutt, and others.
Seventeen come from stenographic notes, one from Persian notes, and one, the
talk at Howard University, is marked as translated by Amin Banani~\citep{PUP}.
The passage on affinity that I read below comes from Esther Foster's notes;
the passage on motion comes from stenographic notes.
\textit{Some Answered Questions}, the source of the third passage, is also a
record of utterances, but of a different kind: \AbdulBaha reviewed its Persian
transcripts himself, and it carries a correspondingly greater weight within
the tradition.

The English of the American talks is therefore the record of an
interpretation, not \AbdulBaha's own wording, and the talks themselves show
why that matters.
Speaking to a federation of women's clubs in Chicago, he pointed out that ``In
Persian and Arabic there are two distinct words translated into English as
man,'' one taking in women and men together and one restricted to the
male~\citep[Talk~32]{PUP}.
A single English word concealed a distinction the originals make.
The English words ``love,'' ``affinity'' and ``attraction'' may conceal
distinctions of the same kind, which the English-speaking listeners had no way
to check.

I draw three working rules from this.
First, I lean on thoughts rather than phrases.
A claim that appears once, in one listener's notes, bears less weight than a
claim that recurs in the notes of many.
The claim that existence consists in affinity passes this test: versions of it
appear in talks taken down by five different listeners in New York,
Washington, Chicago and Eliot, Maine, between April and August
1912~\citep[Talks~2, 23, 41, 48, 74, 89, 92]{PUP}.
Second, where his letters speak to the same point, they take priority.
One of them states the claim outright: ``Through affinity and attraction all
living things like plants, animals and men come into existence, while division
and discord bring about decomposition and
destruction''~\citep[\S 225]{SelectionsAbdul}.
That sentence comes through translators, but not through a note-taker.
Third, no argument in this book rests on a fine shade of English wording,
because that is exactly where the record of an interpretation is least
reliable.

None of this makes the talks untrustworthy as records of what he taught.
The principle I turn to next is a good example of a teaching that stands on
many records rather than one: the harmony of religion and science, to which he
returned with particular frequency.

% ─────────────────────────────────────────────────────────────────────────────
\section{The Two-Wings Principle and What It Predicts}
\label{sec:ch3-two-wings}
% ─────────────────────────────────────────────────────────────────────────────

``Religion and science are the two wings upon which man's intelligence
can soar into the heights, with which the human soul can progress.''

The sentence is from \textit{Paris Talks}~\citeyearpar[p.~143]{ParisTalks},
and it is one of the most frequently quoted statements in \bahai literature.
It is also one of the most frequently misread.

The standard reading treats it as a diplomatic formula --- a generous
acknowledgment that both religion and science have something to contribute,
a plea for mutual tolerance from a tradition trying to establish itself
in a world where religious claims are often at war with scientific ones.
That reading is not wrong, exactly, but it stops well short of what
\AbdulBaha is actually claiming.

Consider what the metaphor requires.
A bird with two wings is not a bird that happens to have two things
attached to it, one of which it uses on Sundays and one on Mondays.
Both wings must work, and both must work in the same direction, and
they must work simultaneously.
A bird that beat only one wing would not merely go more slowly.
It would go in circles.
The two-wings metaphor is a claim about necessary structural coordination
between two methods that, if genuine, must be tracking the same reality.

The deeper claim is this: if religion and science are both functioning
properly --- if each is genuinely tracking the nature of things rather
than merely expressing cultural preferences or institutional self-interest ---
then they should converge.
Not on identical conclusions, arrived at by identical methods with identical
vocabulary.
But on shared structure.
If you and I are both drawing maps of the same territory, using completely
different conventions and surveying from different vantage points, our maps
will look different.
But the underlying structure they describe --- the relationships among
mountains, rivers, and valleys --- should be the same, or our maps are
not both accurate.

By the rules of the previous section, the two-wings sentence, preserved in the
English of those who heard it in Paris, would be a thin foundation on its own.
But the American talks state the principle again and again, in the notes of
different listeners, and in a form the image of the wings leaves out.
In Pittsburgh the record has him saying that ``Any religious belief which is
not conformable with scientific proof and investigation is
superstition''~\citep[Talk~44]{PUP}.
Stated this way, the principle cuts both ways.
It commits the tradition to a test it could fail, and the comparison this book
makes is that test.

Previous dialogue between \bahai thought and science has often worked by
identifying resonances: quantum uncertainty sounds like the unknowability of
God; entanglement sounds like the interconnection of all things; the Big Bang
sounds like the moment of creation.
These resonances are interesting, but they are analogical.
Analogy shows that two things look similar from a certain angle.
Structural isomorphism --- the kind of correspondence I am documenting in
this book --- shows that two frameworks have the same logical form: the same
architecture of claims, the same relationships among their parts, the same
implications for what is possible and what is not.

Analogy is a feature of maps that merely look alike.
Structural isomorphism is what happens when two genuinely accurate maps
of the same territory finally meet.

The question this book pursues is whether the structural correspondence
between the \bahai writings, process philosophy, and quantum mechanics
is of the second, stronger kind.
My argument is that it is.
And the place to begin is with three passages from the writings
that carry most of the philosophical weight.

% ─────────────────────────────────────────────────────────────────────────────
\section{The Three Core Passages}
\label{sec:ch3-passages}
% ─────────────────────────────────────────────────────────────────────────────

These are not the three most famous passages in the \bahai writings.
They are, however, the three that matter most for the argument of this book ---
the three in which a coherent and precise description of the nature of
reality comes most clearly into focus.
I will read each passage closely, more closely than such passages are
usually read outside academic commentary, because their philosophical
content is hidden in plain sight.
The claims are not difficult once you look for them.
But they must be looked for.

\subsection{Existence as Relational Affinity}
\label{subsec:ch3-passage1}

The first passage is from a talk \AbdulBaha gave in New York on 13 May 1912,
at a reception held by the New York Peace Society at the Hotel Astor, and
collected in \textit{The Promulgation of Universal Peace}.
His audience had gathered to promote international peace, and he began with
peace and war: peace builds, and war destroys.
Then he gave a reason for that contrast, and the reason was a claim about
what existence is.

\begin{quote}
``All created things are expressions of the affinity and cohesion of
elementary substances, and nonexistence is the absence of their
attraction and agreement. \ldots\ Everything partakes of this nature and is
subject to this principle, for the creative foundation in all its
degrees and kingdoms is an expression or outcome of love.''
\citep[Talk 48]{PUP}
\end{quote}

Notice first where the passage sits.
It is the premise of an argument for peace, and the talk turns straight from
it to ``the restlessness and agitation of the human world today because of
war''~\citep[Talk~48]{PUP}.
The talk offers the claim as a law of which the human case is one instance:
``Everything partakes of this nature and is subject to this principle.''
Only a premise taken as true can explain why peace builds and war destroys,
and that is my first reason for reading the passage as a claim about the
structure of things.
The same setting counsels caution: the passage was spoken to peace advocates,
not offered as a report on the constitution of matter, and I will not read it
as physics.

Read this slowly.

The first clause makes a claim about what existence \emph{is}.
Not what it is like.
Not what it involves.
What it \emph{is}.
All created things are expressions of the affinity and cohesion of elementary
substances.
Not: all created things are composed of substances that have affinities with
one another.
The affinity and cohesion is not a feature that things have in addition to
existing; it is what their existence consists in.
A thing exists insofar as the relational bonds that constitute it are active.

The second clause makes a claim about what nonexistence is.
Nonexistence is the absence of the attraction and agreement.
This is an unusual claim.
On the ordinary picture, nonexistence is a different kind of state --- the
empty state, the state of absence, the placeholder where a thing is not.
On \AbdulBaha's account, nonexistence is not an independent state at all.
It is what remains when the relational affinity ceases.
You do not go from existence to nonexistence the way you go from a lit room
to a dark room; you go the way a flame goes when the combustion stops, which
is to say, there is no place you go, because there was no flame-substance
that could travel anywhere --- there was only the process of burning, and
now there is not.

The final sentence extends the scope: everything partakes of this nature,
and the creative foundation is at work ``in all its degrees and kingdoms.''
This is not a limited claim about physical matter, later to be qualified
when we get to living things or minds or souls.
The architecture is universal: the same relational principle operates at
every level of reality, from the arrangement of atoms to the orientation
of the soul toward the divine.

And then the same sentence names the principle.
The creative foundation, operating at every degree and in every kingdom, is
``an expression or outcome of love.''

The word this book uses for that principle is \textit{mahabbat}, a Persian
word for love taken from the Arabic \textit{mahabba}; whether it is the word
\AbdulBaha spoke at the Hotel Astor, the English record cannot say, and I
return to the word shortly.
But notice first what is being claimed.
The force that constitutes existence at every level is named love.
Not love as a human emotion first, and then extended by metaphor to the
rest of the universe.
Love as a structural principle that operates universally, whose most
conscious and most intensified form at the human level is what we
recognize as the experience of love.

I am aware that this sounds like a large claim.
Let me be precise about what it does and does not involve.

The claim is not that atoms feel love.
The claim is not that the universe is conscious in the way that persons
are conscious.
The claim is that the relational force that holds an atom together ---
the affinity and cohesion that constitutes its existence --- is the same
structural principle that, at the level of human consciousness, registers
as what we call love.
In the talks' own vocabulary, the difference is one of degree of expression,
though, as I show below, the degrees are separated by real thresholds.
What operates at the level of the proton as attraction and cohesion
operates at the level of the human soul as longing, devotion, and care.

This is a philosophically precise claim with specific logical consequences.
It means that the question ``why does love exist?'' has a different shape
than it does on the usual view.
On the usual view, love is a product of biological evolution in beings
complex enough to have social lives, explained by the pressures that selected
for pair-bonding, parental investment, and group solidarity.
That explanation may be accurate as far as it goes.
But on \AbdulBaha's account, it does not go deep enough.
Evolution did not produce a new thing called love in organisms complex
enough to feel it.
Evolution produced organisms complex enough to consciously experience the
relational principle that was already operating at every level of the physical
world.
The love came first.
The organisms came along later and became capable of noticing it.

Now look at the structure of the claim from the angle of physics.

\citet{Bell1964} proved, in a landmark 1964 paper, that the correlations
exhibited by quantum systems cannot be explained by any theory in which
distant particles have pre-existing, locally determined properties.
If no influence travels faster than light, the correlations cannot be the
result of particles already having their properties fixed and then being
measured.
On the relational reading of this result that I defend in this book, and that
not every physicist shares, the properties come into being through the
measurement event --- through the relational interaction between particle and
measuring apparatus.
Properties do not precede the relations that manifest them;
properties are constituted through those relations.

This is the same logical structure as \AbdulBaha's first passage.
Things do not first exist and then enter into affinities with one another.
The affinities constitute the things.
Existence is the relational event, not its precondition.

I am not claiming that \AbdulBaha was doing quantum mechanics.
I am claiming that the logical form of his ontological claim and the logical
form of the conclusion that, on this reading, physics draws from Bell's
theorem and from the experiments honored by the 2022 Nobel Prize in
Physics~\citep{Aspect1982} are the same.
Two independent traditions, using entirely different methods, arrived at
the same structural description of the elementary level of reality.
That is what I mean by structural isomorphism.
That is what the two-wings principle, taken seriously, predicts.

\subsection{Motion and the Primacy of Process}
\label{subsec:ch3-passage2}

The second passage opens a talk given nine days after the first, on 22 May
1912, to a Unitarian conference at the Tremont Temple in Boston.
These are its first words:

\begin{quote}
``Creation is the expression of motion. Motion is life. A moving object
is a living object, whereas that which is motionless and inert is as dead.
All created forms are progressive in their planes, or kingdoms of
existence, under the stimulus of the power or spirit of life. The
universal energy is dynamic. Nothing is stationary in the material world
of outer phenomena or in the inner world of intellect and consciousness.''
\citep[Talk 52]{PUP}
\end{quote}

This passage states one of the oldest and most contested ideas in
philosophy in a form that is, once you see it, unusually clear.

The dominant assumption of Western philosophy for more than two thousand
years --- from Aristotle through Newton and into the twentieth century ---
is that the world is fundamentally composed of things.
Things have properties and enter into relations.
Motion is something that happens to things; rest is their natural state.
The billiard ball sits on the table until something moves it; the planet
travels in its orbit, but what it really is is the planet, and the orbit
is what it does.

\AbdulBaha reverses this.
Creation is the expression of motion.
Motion is life.
What is motionless and inert is as dead.

Notice the grammar of the first two sentences.
He does not say that creation involves motion, or that creation produces
motion, or that motion is a sign of life.
He says that creation \emph{is} the expression of motion.
He says that motion \emph{is} life.
These are equations of essence, not observations about correlation.
He is not saying that living things tend to move.
He is saying that motion --- dynamic relational process --- is what existence
\emph{is}, and that what we call ``things'' are patterns of motion that have
achieved sufficient regularity to appear stable.

The talk then puts this claim to work, as the Peace Society address did.
Its next paragraph draws the consequence the speaker cared about: ``Religion
is the outer expression of the divine reality. Therefore, it must be living,
vitalized, moving and progressive''~\citep[Talk~52]{PUP}.
A religion without motion, the talk goes on, is dead.
Once again a general law about existence is applied to human life.

The stone on the table looks nothing like motion.
But go far enough into the stone and you find a seething world of atomic
and subatomic activity: electrons bound to nuclei, quarks bound together
by the exchange of gluons, the whole structure sustained by relational
processes operating at energies our macroscopic intuition cannot register.
The stone is not a thing that happens to have some internal motion.
The stone is the pattern of all that relational activity, perceived
at a scale and a speed that makes the pattern appear static.

Physics has confirmed this in a form more rigorous than any metaphysical
intuition can match.
\citeauthor{Heisenberg1927}'s uncertainty principle, formulated in 1927,
proves that perfect rest is not an extreme case of limited motion;
it is physically impossible~\citeyearpar{Heisenberg1927}.
You cannot specify a particle's position and its momentum simultaneously
with arbitrary precision, because position and momentum are not
independently real quantities that a particle merely happens to have.
They are aspects of the relational process by which the particle's state
is constituted through interaction with measuring devices.
The vacuum itself --- what we might think of as the ultimate motionless
void --- is not empty and still.
It seethes with quantum fluctuations: the residual agitation of quantum
fields that no removal of energy can eliminate.
Perfect inertness, as \AbdulBaha suggests, is not merely unusual; it is
the absence of physical reality itself.

Here again the logical structure of the \bahai claim and the logical
structure of the physics conclusion coincide.
In both cases, the verdict is: reality is dynamic through and through;
the appearance of static substance is a macroscopic averaging of
relational process; and what we call a thing is always, at bottom,
a pattern of becoming.

This is the convergence with process philosophy that I will develop
in Part Two.
For now, let me note the connection and leave it there: Alfred North
Whitehead, writing in 1929, would build an entire philosophical
framework on the claim that events are primary and things are derivative
--- that the fundamental unit of reality is not the persistent object
but the momentary occasion of becoming~\citeyearpar{Whitehead1929}.
The \bahai writings had stated the conclusion of that framework decades
before the framework existed.

\subsection{The Soul and the Mirror}
\label{subsec:ch3-passage3}

The third passage is different in character from the first two.
The first two passages address the structure of physical reality.
This one addresses the nature of the soul, and specifically its relationship
to the physical body.
It comes from \textit{Some Answered Questions}, a record of \AbdulBaha's
answers to questions put to him at table in `Akká between 1904 and 1906,
while he was still a prisoner there; the American \bahai Laura Clifford
Barney, who asked many of the questions, compiled the book.

\AbdulBaha writes there that the soul ``has a connection with the body like
that of the sun with this mirror''~\citep[Ch.~68]{SAQ}. The sun is not inside
the mirror, yet the two are connected, and on his account what befalls the
mirror does not befall the sun.

The image is vivid, but its philosophical content is easy to miss unless
you hold it against the claims of the preceding two passages.

If existence is constituted by relational affinity (passage one) and if
reality is fundamentally dynamic process (passage two), then the soul must
be described within that framework.
The question is: what kind of entity is the soul within a relational,
process ontology?

\AbdulBaha's answer is that the soul is a non-composite entity.
It is not made of parts in the way that a body is made of parts.
It does not have constituent elements whose relational affinity constitutes
its existence in the way that a stone's existence is constituted by the
relational affinity of its component atoms.
In the words of the American talks, the human spirit is ``simple, single and
not composed''~\citep[Talk~100]{PUP}; it ``is not a composition of
elements''~\citep[Talk~85]{PUP}.
\AbdulBaha makes the same point in \textit{Some Answered Questions} and
\textit{Paris Talks}, and the consequence is significant: because the soul does not exist in
virtue of relational affinities that could come apart, it is not subject
to the dissolution that awaits everything whose existence consists in
such affinities.

The sun-and-mirror image captures this with precision.
The mirror is a physical object: composite, material, existing in virtue
of the arrangement of its molecules, subject to breakage.
When the mirror breaks, the arrangement dissolves, the form perishes,
the object passes from existence.
The sun is a different kind of entity: it illuminates the mirror, it is
connected to the mirror in the sense that the mirror's function depends
on the sun's presence, but the sun is not \emph{in} the mirror and is
not affected by the mirror's destruction.

The soul, on this account, is connected to the body in a way that makes
possible the body's highest functions --- perception, thought, moral
awareness, love --- while being itself of a different ontological order.
The body is the local apparatus through which the soul's reality becomes
manifest in the physical world.
The soul is the non-local source of that manifestation.

This is not Cartesian dualism, though the surface similarity is easy to
mistake.
Descartes' dualism posited two separate substances --- thinking substance
and extended substance --- and then faced the insoluble problem of how they
could interact.
\AbdulBaha is not positing two substances.
He is positing a relational architecture: the soul is the source,
the body is the apparatus of manifestation, and the connection between
them is functional and relational, not substance-based.
When the body ceases --- when its constituent affinities dissolve ---
the connection ceases, but the source does not.

What this passage is doing, in relation to the first two, is extending
the relational framework into its most philosophically consequential
territory.
This third passage describes what lies at the boundary of that physical
existence: an entity that participates in the physical world through
relational connection but is not itself constituted by the kind of
relational affinity that makes physical things the things they are.

The sun does not become less the sun when this particular mirror
happens to be in its light.
The soul, on this account, does not become less itself because it happens,
in this particular life, to be connected to this particular body.

I am going to return to this passage later in the book --- in
Chapter~\ref{ch:paradigm-required}, where I develop what it might mean for
the soul to be a non-composite reality within a framework that has been
clarified by both process philosophy and quantum mechanics.
For now, the point I want to establish is simpler: this is an ontological
claim, not a poetic one.
Its logical form is precise.
It says something specific about the kind of entity the soul is, and
that something specific fits within the coherent metaphysical framework
that the other two passages have begun to build.

% ─────────────────────────────────────────────────────────────────────────────
\section{What the Writings Claim About Reality}
\label{sec:ch3-claims}
% ─────────────────────────────────────────────────────────────────────────────

I want to step back from the individual passages and see what they
add up to.

Three claims have been made.

The first is a claim about existence: existence at every level, in every
kingdom, consists in the relational affinity of constituents.
What is not in relation is not in existence.
Nonexistence is the cessation of the relational bond, not a separate state.
The force sustaining the relational bond is named love; this book's word for
it is \textit{mahabbat}.

The second is a claim about the fundamental character of reality:
reality is dynamic, processual, constituted by motion and becoming.
What appears static is a pattern of relational process averaging out over
a scale at which individual events are not perceptible.
Perfect inertness is identical with non-being.

The third is a claim about the soul: the soul is a non-composite entity
connected to the physical body through a functional relational architecture
but not constituted by the kind of relational affinity that makes physical
things the things they are.
The soul is the source; the body is the apparatus of manifestation.
Connection, not composition.

Together, these three claims constitute a coherent ontology.
It is not a complete philosophical system --- the writings are not a
treatise in systematic metaphysics --- but it is a set of structural claims
with a definite logical form, and that form is identifiable.

The logical form can be stated in one sentence: \emph{existence is
constituted by relational process; process is primary, substance is derivative;
and love --- \textit{mahabbat} --- is the name for the force that constitutes
relational existence at every degree.}

This is not a poetic description.
It is not a metaphor for something that could be said better in other terms.
It is a philosophical claim about the structure of reality --- one that was
made in the period between the 1850s and 1921, in the context of what the
\bahai writings call revelation and philosophical reflection, and that can
now be evaluated against what physics and philosophy have independently arrived
at in the century and a half since.

\subsection{The Word \textit{Mahabbat} and What It Names}

Before moving on, I want to linger on the word that does the most work
in this framework.

\textit{Mahabbat} is the Persian form of an Arabic word, and it belongs to a
family of terms for love.
It shares a root with \textit{hubb}, the most general of them, and it stands
beside \textit{mawaddat} (affectionate kindness, often interpersonal) and
\textit{`ishq} (the ardent, consuming love of the Sufi mystical tradition).
These distinctions are not pedantic.
They show that the Persian and Arabic vocabulary available to the writers of
the \bahai texts had finer tools for discussing love than English does, and
that one English word can stand for several of them.

For the same reason I have to be careful: the English records of the talks do
not say which word stood behind each ``love,'' and I do not pretend that they
do.
My reason for choosing \textit{mahabbat} as this book's key term lies in its
sense, not in the wording of any single talk.
\textit{Mahabbat} carries the semantic weight of cordial inclination as a
sustained relational disposition.
It is the kind of love that draws toward, that tends toward, that constitutes
the sustained orientation of one reality toward another.
That is the sense the principle needs, and it matches the description of love
in one of \AbdulBaha's letters as ``the unique power that bindeth together the
divers elements of this material world''~\citep[\S 12]{SelectionsAbdul}.

The argument this book makes is that the principle I am calling
\textit{mahabbat} is doing exactly the same structural work as \textit{eros}
in Whitehead --- Whitehead's term for the creative drive toward richer, more complex relational
integration that he identifies as operative at every level of reality.
Both are names for the universe's intrinsic tendency toward relational
coherence.
Both operate at every scale.
Both are not additions to the basic mechanics of the world but names for
the world's basic directional tendency.

When I say that the human experience of love is the most conscious form of
\textit{mahabbat}, I am not diminishing love by explaining it away.
I am doing the opposite.
I am saying that what you feel when you love someone is not a subjective overlay
on a universe that would be exactly the same without you --- a universe of
particles colliding according to laws that have nothing to do with care,
beauty, or belonging.
What you feel when you love someone is your most direct access to the structural
principle that constitutes reality at every level.
The love is not in you, added to an otherwise loveless world.
The love is in the world, and you are, among all the entities in the world,
one of the ones who has become conscious enough to feel it.

This is the ontological claim the three passages put on the table.
The rest of this book examines whether physics and philosophy give us reasons
to take that claim seriously.

\subsection{The Scope Is Universal}

``Everything partakes of this nature \ldots in all its degrees and kingdoms.''

The phrase ``all its degrees and kingdoms'' is not rhetorical amplification.
It is a precise scope claim.
The \bahai writings distinguish several ``kingdoms'' or levels of
reality: the mineral, the vegetable, the animal, the human, and what they
sometimes call the divine.
To say that the relational principle of \textit{mahabbat} operates in all its
degrees and kingdoms is to say that the same architectural principle governs
every level --- not that love is present everywhere in the same way, but that
the fundamental relational structure is the same at every level.

In August 1912, at Green Acre in Eliot, Maine, \AbdulBaha gave two talks on
consecutive days, both taken down by Edna McKinney, that climb the kingdoms
one rung at a time.
The first sets out to show that ``love is the cause of the existence of all
phenomena and that the absence of love is the cause of disintegration or
nonexistence.''
In the mineral, it says, the power of cohesion ``is in reality love or
affinity manifested in a low degree according to the exigencies of the
mineral world''; in the plant, attraction brings growth; in the animal, it
adds feeling and fellowship; and in the human being there is, beyond all
these, ``the attraction of heart, the susceptibilities and affinities which
bind men together''~\citep[Talk~89]{PUP}.
The second talk puts the lowest rung in a phrase I find exact: the power of
attraction in the stone is love, ``the only expression of love the stone can
manifest''~\citep[Talk~92]{PUP}.
To say that love holds the stone together, then, says nothing about a stone's
feelings; it says that cohesion is what love amounts to at a stone's capacity.

Two features of these talks keep the reading honest.
The rungs are not a smooth slope.
The second talk calls the human degree of attraction, which it describes as
``conscious and spiritual,'' ``an immeasurable advance''~\citep[Talk~92]{PUP},
and an earlier talk insists that ``A lower degree cannot comprehend a
higher''~\citep[Talk~46]{PUP}, a claim Chapter~\ref{ch:observer-dependence}
takes up.
And the first talk does not end with human love.
Having reached the human kingdom, it steps back: ``But these are only degrees
of love which exist in the natural or physical world.''
``Real love,'' it continues, ``is the love which exists between God and His
servants, the love which binds together holy souls''~\citep[Talk~89]{PUP}.
On the talk's own terms, the attraction that binds the elements and the
affection that binds people are true degrees of love, but not its highest
form.
A reading of this book that reduced the love of God to a physical force would
have that talk against it, and me as well.

This has a consequence that matters for the argument of this book.
If \textit{mahabbat} is the creative foundation of existence at the level of
elementary substances, and the same principle operates at the level of
human consciousness, then the difference between a rock and a person is not
the difference between a thing that has no love and a thing that has love.
It is the difference between a pattern of relational activity so simple
that no experience of it is registered, and a pattern of relational activity
so complex that the pattern becomes capable of experiencing itself.

On the process reading I develop in Part Two, consciousness is not a new kind
of substance added to the world at some point in evolutionary history.
It is what the relational structure of reality looks like from the inside when
the structure has achieved sufficient complexity.
The universe was always doing what conscious beings experience.
At some point in its development, it became complex enough to notice.

The \bahai writings assert the continuity of love across the kingdoms, but
they do not say that the human soul is produced by complexity.
As the third passage showed, the soul on their account is of a different order
from the body, which is the mirror in which it appears and not its source.
Chapter~\ref{ch:process-ontology} gives the process account of the soul as a
pattern, and Chapter~\ref{ch:paradigm-required} takes up the \bahai claim that
the soul is not composite; they are not the same claim, and I do not treat
them as one.

The process-philosophical framework of Part Two gives the continuity a more
rigorous form, and the physics of Part One supplies the relational floor on
which it stands, though physics alone says nothing about experience.
For now, the point is that it is a claim --- a precise, structural claim
with logical consequences --- and one the \bahai writings make with full
seriousness and evident care.

\subsection{On the Chronological Sequence}

The passages I have analyzed were delivered between 1904 and 1912.
The underlying writings of \Bahaullah from which \AbdulBaha draws go back
to the 1850s.
Alfred North Whitehead published \textit{Process and Reality}, which gives
these ideas their most rigorous philosophical articulation in the Western
tradition, in 1929~\citeyearpar{Whitehead1929}.
John Bell published the theorem that rules out classical substance ontology,
in its local form, as a description of quantum phenomena in 1964~\citeyearpar{Bell1964}.
The experimental confirmations of Bell's theorem that earned the 2022 Nobel
Prize in Physics were conducted decades after that~\citeyearpar{Aspect1982}.

I draw attention to this sequence now and will return to it in
Chapter~\ref{ch:bahai-prior}.
The point I want to make here is negative: it establishes independence.
The \bahai writings were not borrowing from Whitehead.
\AbdulBaha was not reading physics papers.
The structural convergence is not the result of influence in any direction.

What the chronological sequence establishes is this: whatever the source of
the \bahai writings' philosophical conclusions, those conclusions were arrived
at by a different method, in a different language, addressed to a different
audience, from a different civilizational context, before the other two
traditions existed in their modern forms.
The convergence, when it becomes apparent, is the convergence of independent
investigations of the same reality.
That is what the two-wings principle predicts.

\subsection*{The Charge of Anachronism}

The chronology invites the strongest objection to this chapter.
The talks were given in 1912, in the physical vocabulary of their day, to
audiences who had never heard of entanglement.
Reading quantum mechanics into them, the objection runs, credits the texts
with ideas that did not yet exist, by picking out the sentences that happen to
sound modern.

Much of this is true, and the talks supply the evidence.
One has the elements of a tree, a man or a fish combining ``in chemical
affinity,'' and a mirror, a table or a clock held together by ``a cohesive
power''~\citep[Talk~74]{PUP}.
Another explains sight by ``vibration or movement in the ethereal
element''~\citep[Talk~58]{PUP}, the luminiferous ether that physics was then
in the course of abandoning.
Anyone who says that these talks describe entanglement is making the mistake
the objection names, and where their physics is dated, no reading rescues the
detail.

The objection does not reach the claim I make, for three reasons.
I compare logical forms, and the form of the Peace Society claim --- that a
composite thing exists in the binding of its constituents and ceases when the
binding fails --- can be stated in the vocabulary of 1912 or of 2022 without
crediting the talk with any idea of an entangled state.
The charge of selective quotation can be checked, and the affinity claim
survives the check: it recurs in the notes of five listeners and in a letter
from \AbdulBaha's own pen.
And the texts invite the comparison, since the principle of the two wings asks
religious claims to answer to science, the science of later ages included.
The correspondence I claim lies at the level of structure, and only there.

% ─────────────────────────────────────────────────────────────────────────────
\section{Confirmation, Not Appropriation}
\label{sec:ch3-confirmation}
% ─────────────────────────────────────────────────────────────────────────────

A clarification is essential, and I want to make it before the reader
reaches the chapters where the detailed convergences are documented.

I do not claim that the \bahai writings are physics texts.
I do not claim that \Bahaullah or \AbdulBaha was doing quantum mechanics
or process philosophy under other names.
I do not claim that the writings carry hidden physical discoveries awaiting
extraction by sufficiently clever readers.
I do not claim that the English of the talks is \AbdulBaha's own wording.
Nor do I claim that my reading is the \bahai community's: authoritative
interpretation in that tradition belongs to \AbdulBaha and to Shoghi Effendi,
whom he appointed to succeed him, and mine is one scholar's reading, which a
\bahai reader is free to find wanting.

The \bahai writings are spiritual writings, addressed to questions of human
meaning, moral development, the nature of God, and the obligations of
human beings toward one another.
They are not seeking to describe the behavior of quantum systems or to
adjudicate among interpretations of wave function collapse.
Their concerns are different from those of physics, and their methods
are different, and their vocabulary is different.

What I claim is different from all of that.

I claim that the \emph{structural description} of reality in the \bahai writings
--- existence as relational affinity, reality as dynamic process, the force
constituting relational existence named love --- corresponds to the structural
description that quantum mechanics has independently arrived at through
a century of mathematics and experiment, and that Whitehead's process
philosophy has independently arrived at through philosophical analysis.

Each tradition keeps its own vocabulary.
Each tradition validates its conclusions on its own grounds.
A \bahai scholar does not need to accept the formalism of quantum field
theory to validate the relational reading of the writings.
A physicist does not need to accept the authority of the \bahai revelation
to accept the relational interpretation of quantum mechanics.
Each tradition, on its own terms, arrives at the same structural conclusions.
The convergence confirms that the structure is real --- that it is not an
artifact of any one tradition's particular vocabulary or assumptions.

The relationship among the three traditions is confirmation, not
appropriation~\citeyearpar{AdamsMahabbat}.
The \bahai writings do not need quantum mechanics to be true.
Quantum mechanics does not need the \bahai writings to be true.
What I document is that, independently of each other, all three arrived at
the same description of reality --- which is precisely what you would expect
if all three were genuinely tracking the same thing.

The two-wings metaphor, taken at full strength, predicts exactly this.
If religion and science are both genuinely tracking reality, they should
confirm each other's conclusions at the level of structure.
Not because one borrowed from the other.
Because the structure is there to be found, and genuine investigation
finds it.

\begin{convergencemoment}{The Writings: \textit{Mahabbat}, the Two Wings, and the Nature of Reality}
Three traditions, using three different methods, arrived at the same
structural description of reality.

From the quantum side: Bell's theorem, confirmed by the Nobel Prize-winning
experiments, proves that properties do not precede the relational interactions
that manifest them~\citeyearpar{Bell1964,Aspect1982}.
Existence, at the elementary level, is constituted by relational interaction.
There is no independent, self-contained substance beneath the relations.

From the process-philosophical side: Whitehead's actual occasions are not
things but relational events of becoming; what we call ``things'' are societies
of such occasions maintaining patterns of relational process;
the fundamental unit of reality is the event, not the object~\citeyearpar{Whitehead1929}.

From the \bahai side: ``All created things are expressions of the affinity
and cohesion of elementary substances, and nonexistence is the absence of
their attraction and agreement \ldots\ for the creative foundation in all its
degrees and kingdoms is an expression or outcome of love''~\citep[Talk~48]{PUP}.

One sentence: existence is relational becoming --- and the name for the force
that constitutes that becoming, at every degree and in every kingdom, is love.
\end{convergencemoment}
